\begin{table}[htbp]
\centering
\caption{Sample composition by SIC major industry division under the primary specification: firm-years, distressed observations, and distress rate (\%), FY1990--2023. Distress is a CRSP performance delisting, excluding voluntary and administrative exits, within the 12-month window following the 10-K filing date. Financials (SIC 6000--6999) and utilities (4900--4999) are excluded from the primary sample.}
\label{tab:sample_sic}
\footnotesize
\setlength{\tabcolsep}{4pt}
\begin{tabular}{lrrr}
\toprule
sic\_div & n\_total & n\_distressed & distress\_rate\_pct \\
\midrule
A: Agriculture (0–999) & 528 & 9.0000 & 1.7050 \\
B: Mining (1000–1499) & 5341 & 151.0000 & 2.8270 \\
C: Construction (1500–1999) & 1577 & 49.0000 & 3.1070 \\
D: Manufacturing (2000–3999) & 58247 & 1241.0000 & 2.1310 \\
E: Transp/Comm (4000–4999, util. excluded) & 6335 & 172.0000 & 2.7150 \\
F: Wholesale (5000–5199) & 4465 & 125.0000 & 2.8000 \\
G: Retail (5200–5999) & 8702 & 245.0000 & 2.8150 \\
H: Services I (7000–7999) & 18578 & 492.0000 & 2.6480 \\
I: Services II (8000–8999) & 5800 & 163.0000 & 2.8100 \\
J: Other (9000–9999) & 1264 & 113.0000 & 8.9400 \\
\bottomrule
\end{tabular}
\end{table}
